%%============================================================================
%% MANIFESTO — Grupo de Usuários Linux e Software Livre - Estatística e Ciência de Dados
%%             Universidade Federal do Paraná (UFPR)
%%============================================================================
%% Documento institucional — versão formal
%% Compilar com: pdflatex manifesto.tex  (ou xelatex / lualatex)
%%============================================================================

\documentclass[
  12pt,          % tamanho base da fonte
  a4paper,       % formato de papel
  oneside,       % impressão frente única
  brazil         % idioma principal: português do Brasil
]{article}

%---------------------------------------------------------------------------
% Pacotes essenciais
%---------------------------------------------------------------------------
\usepackage[T1]{fontenc}            % codificação de fonte
\usepackage[utf8]{inputenc}         % entrada UTF-8 (pdflatex)
\usepackage{lmodern}                % fontes Latin Modern (vetoriais)
\usepackage{microtype}              % microtipografia (melhora legibilidade)

\usepackage[brazil]{babel}          % hifenização e nomes de seções em PT-BR

\usepackage[
  margin=2.5cm,
  top=3cm,
  bottom=3cm,
  headheight=25pt
]{geometry}                         % margens do documento

\usepackage{setspace}               % controle de espaçamento entre linhas
\usepackage{xcolor}                 % suporte a cores
\usepackage{titlesec}               % formatação de títulos de seções
\usepackage{titling}                % formatação do título do documento
\usepackage{fancyhdr}               % cabeçalhos e rodapés personalizados
\usepackage{enumitem}               % personalização de listas
\usepackage{lettrine}               % capitulares (letra inicial decorada)
\usepackage{graphicx}               % suporte a imagens (se necessário)
\usepackage{hyperref}               % hiperlinks e metadados do PDF
\usepackage{bookmark}               % bookmarks do PDF mais limpos

%---------------------------------------------------------------------------
% Cores institucionais
%---------------------------------------------------------------------------
\definecolor{gulecdBlue}{HTML}{1B3A5C}   % azul escuro institucional
\definecolor{gulecdAccent}{HTML}{2E86AB} % azul de destaque
\definecolor{gulecdGray}{HTML}{555555}   % cinza para texto secundário
\definecolor{gulecdLight}{HTML}{F5F7FA}  % fundo claro para caixas
\definecolor{gulecdRule}{HTML}{D0D5DD}   % cor de filetes e réguas

%---------------------------------------------------------------------------
% Metadados do PDF
%---------------------------------------------------------------------------
\hypersetup{
  pdftitle    = {Manifesto — GULECD/UFPR},
  pdfauthor   = {Grupo de Usuários Linux em Estatística e Ciência de Dados},
  pdfsubject  = {Manifesto institucional do GULECD},
  pdfkeywords = {software livre, código aberto, ciência de dados, UFPR, manifesto},
  colorlinks  = true,
  linkcolor   = gulecdAccent,
  urlcolor    = gulecdAccent,
  citecolor   = gulecdAccent
}

%---------------------------------------------------------------------------
% Configurações tipográficas
%---------------------------------------------------------------------------
\onehalfspacing                     % espaçamento entre linhas = 1,5
\setlength{\parindent}{1.5cm}       % recuo de parágrafo
\setlength{\parskip}{0.5em}         % espaço entre parágrafos

%---------------------------------------------------------------------------
% Formato do título do documento
%---------------------------------------------------------------------------
\pretitle{%
  \begin{center}
  \vspace*{1cm}
  \rule{\textwidth}{2pt}\vspace{0.3cm}\\
  \LARGE\bfseries\color{gulecdBlue}
}
\posttitle{%
  \\\vspace{0.3cm}
  \rule{\textwidth}{2pt}
  \end{center}
  \vspace{0.5cm}
}
\preauthor{\begin{center}\large\color{gulecdGray}}
\postauthor{\end{center}}
\predate{\begin{center}\small\color{gulecdGray}}
\postdate{\end{center}}

%---------------------------------------------------------------------------
% Formato das seções
%---------------------------------------------------------------------------
\titleformat{\section}
  [block]
  {\normalfont\Large\bfseries\color{gulecdBlue}}
  {\thesection.}{0.8em}{}
  [\vspace{0.2em}{\color{gulecdRule}\rule{\linewidth}{0.6pt}}]

\titleformat{\subsection}
  [block]
  {\normalfont\large\bfseries\color{gulecdAccent}}
  {\thesubsection.}{0.8em}{}

\titlespacing*{\section}
  {0pt}{2.5ex plus 1ex minus .2ex}{1.5ex plus .2ex}
\titlespacing*{\subsection}
  {0pt}{2ex plus 1ex minus .2ex}{1ex plus .2ex}

%---------------------------------------------------------------------------
% Cabeçalho e rodapé
%---------------------------------------------------------------------------
\pagestyle{fancy}
\fancyhf{}
\fancyhead[L]{\small\color{gulecdGray}\textsc{GULECD/UFPR}}
\fancyhead[R]{\small\color{gulecdGray}\textit{Manifesto}}
\fancyfoot[C]{\small\color{gulecdGray}---\ \thepage\ ---}
\renewcommand{\headrulewidth}{0.4pt}
\renewcommand{\headrule}{%
  \hbox to\headwidth{%
    \color{gulecdRule}\leaders\hrule height \headrulewidth\hfill}}

% Remove cabeçalho/rodapé da primeira página (será redefinido manualmente)
\fancypagestyle{plain}{%
  \fancyhf{}%
  \fancyfoot[C]{\small\color{gulecdGray}---\ \thepage\ ---}%
  \renewcommand{\headrulewidth}{0pt}%
}

%---------------------------------------------------------------------------
% Ambiente personalizado para os princípios
%---------------------------------------------------------------------------
\newenvironment{principio}[1]{%
  \vspace{0.8em}
  \noindent
  \begin{minipage}{\textwidth}
    \subsection*{#1}
    \itshape
}{%
  \end{minipage}
  \vspace{0.5em}
}

%---------------------------------------------------------------------------
% Informações do documento
%---------------------------------------------------------------------------
\title{Manifesto do Grupo de Usuários Linux e Software Livre\\Estatística e Ciência de Dados - UFPR}
%\author{Grupo de Usuários Linux e Código Aberto\\Curso de Estatística e Ciência de Dados --- UFPR}
\date{\today}

%============================================================================
% INÍCIO DO DOCUMENTO
%============================================================================
\begin{document}

%--- Capa / Título principal -------------------------------------------------
\thispagestyle{empty}
\maketitle

\newpage

%--- Preambulo --------------------------------------------------------------
\section*{Preâmbulo}
\addcontentsline{toc}{section}{Preâmbulo}

\lettrine[lines=3, lhang=0.1, nindent=0em, slope=-0.5em]{E}{ste} grupo nasce da convergência entre a tradição científica da Estatística e da 
Ciência de Dados e a filosofia do \textit{software} livre que sustenta o ecossistema Linux.

Reconhecemos que a produção de conhecimento rigoroso, verificável e cumulativo depende de transparência metodológica, acesso irrestrito a 
ferramentas e capacidade de auditoria independente.

Em um contexto onde infraestruturas proprietárias frequentemente limitam a reprodutibilidade e obscurecem processos analíticos, 
posicionamo-nos deliberadamente a favor de práticas abertas, colaborativas e tecnicamente fundamentadas.

Este manifesto estabelece os princípios que orientam nossas atividades, consolidando um compromisso coletivo com excelência científica, 
autonomia tecnológica e construção compartilhada do conhecimento.

%--- Princípios -------------------------------------------------------------
\section*{Princípios}
\addcontentsline{toc}{section}{Princípios}

\begin{principio}{I --- Primazia do conhecimento aberto}
Defendemos que o conhecimento científico deve ser livremente acessível, auditável e reutilizável; promovemos práticas que favoreçam 
transparência, reprodutibilidade e disseminação irrestrita do saber.
\end{principio}

\begin{principio}{II --- Software livre como fundamento científico}
Adotamos e incentivamos o uso de software livre e de código aberto como base metodológica, reconhecendo sua superioridade epistemológica na 
validação, inspeção e evolução contínua de ferramentas analíticas.
\end{principio}

\begin{principio}{III --- Reprodutibilidade como princípio inegociável}
Toda análise deve ser passível de reprodução independente; priorizamos ambientes versionados, código documentado e \textit{pipelines} de análise de dados auditáveis
 como requisitos mínimos de rigor científico.
\end{principio}

\begin{principio}{IV --- Autonomia tecnológica}
Valorizamos a independência em relação a plataformas proprietárias, promovendo o domínio técnico sobre sistemas, linguagens e ferramentas que 
permitam controle integral do processo computacional.
\end{principio}

\begin{principio}{V --- Cultura de colaboração radical}
Incentivamos a colaboração aberta, horizontal e meritocrática; compartilhamos código, dados e ideias com o objetivo de acelerar a produção científica coletiva.
\end{principio}

\begin{principio}{VI --- Aprendizado contínuo e profundo}
Comprometemo-nos com a formação técnica rigorosa, buscando compreender não apenas o uso, mas também os fundamentos computacionais, 
estatísticos e matemáticos das ferramentas que utilizamos.
\end{principio}

\begin{principio}{VII --- Transparência e rastreabilidade dos dados}
Defendemos que dados e transformações analíticas sejam plenamente rastreáveis, com documentação clara de origem, processamento e limitações, 
garantindo integridade científica.
\end{principio}

\begin{principio}{VIII --- Infraestrutura aberta e escalável}
Priorizamos o uso de ambientes Linux e ferramentas abertas para construção de \textit{softwares} e \textit{pipelines} analíticos robustos, escaláveis e eficientes, 
desde ambientes locais até sistemas distribuídos.
\end{principio}

\begin{principio}{IX --- Ética no uso e compartilhamento de dados}
Aliamos abertura com responsabilidade; respeitamos privacidade, legislação e princípios éticos, garantindo que a liberdade tecnológica não 
comprometa direitos individuais.
\end{principio}

\begin{principio}{X --- Compromisso com impacto científico e social}
Utilizamos ciência de dados e \textit{software} livre como instrumentos para gerar conhecimento relevante, promover inovação e 
contribuir de forma concreta para a sociedade.
\end{principio}

%--- Encerramento -----------------------------------------------------------
\section*{Encerramento}
\addcontentsline{toc}{section}{Encerramento}

Este compromisso não se encerra em princípios abstratos, mas se materializa nas práticas cotidianas de estudo, pesquisa, desenvolvimento e 
colaboração que sustentam este grupo.

Ao adotar e promover \textit{software} livre e código aberto como pilares da Estatística e Ciência de Dados, assumimos uma postura ativa na construção de um ecossistema 
científico mais transparente, robusto e acessível.

Subscrevemos, os membros do Grupo de Usuários Linux e Software Livre do Curso de Estatística e Ciência de Dados da UFPR, 
este manifesto como expressão formal de nossos valores, orientando nossas ações presentes e futuras.

\vspace{2cm}

%============================================================================
\end{document}
%============================================================================
